\chapter{Asset Managers, Pensions, Sovereign Funds and Insurers}\label{ch:in:asset-managers-pensions-sovereign-funds-and-insurers}

Norway's sovereign fund was worth 21\,268 billion kroner at the end of 2025, about \$2.05 trillion at the year's average
exchange rate, and its manager employed 678 people. It spent 3.8 basis points of the fund on management, and 39\% of that
went to the external managers who ran 5\% of its assets. Canada's national pension fund, less than a third of the size,
paid its external managers C\$4.7 billion in base and performance fees in its latest year, 2.7 times its own
operating expenses. The two institutions have made opposite choices about what to do themselves, and those choices decide how
many quants they employ and what the quants do. This chapter describes the other side of the markets: the \omterm{def:in:asset-managers-pensions-sovereign-funds-and-insurers:owner}{asset owners}
whose money the rest of the industry manages, the asset managers who serve them, and the insurers who manage against
their liabilities.

\section{Asset owners and their in-house teams}

\begin{definition}[Asset owner]\label{def:in:asset-managers-pensions-sovereign-funds-and-insurers:owner}
An \emph{asset owner}\index{asset owner} is an institution that invests assets it holds on behalf of beneficiaries or of
the state -- a pension fund, a \omterm{def:in:asset-managers-pensions-sovereign-funds-and-insurers:swf}{sovereign wealth fund}, an endowment, an insurer's general account -- and answers for the
results to them, whether it manages the assets itself or through external managers.
\end{definition}

\begin{definition}[Sovereign wealth fund, defined-benefit pension plan]\label{def:in:asset-managers-pensions-sovereign-funds-and-insurers:swf}
A \emph{sovereign wealth fund}\index{sovereign wealth fund} is a special-purpose investment fund owned by a general
government, created for macroeconomic purposes, that invests (among other things) in foreign financial assets. A
\emph{defined-benefit pension plan}\index{defined-benefit pension plan} promises its members a fixed, pre-established
benefit at retirement, so that the investment risk falls on the sponsor and the plan's assets are managed against the
promised benefits.
\end{definition}

Book~1 (chapter~3) introduced \omterm{def:in:asset-managers-pensions-sovereign-funds-and-insurers:owner}{asset owners} as the clients of asset managers. As employers they differ from everyone else
in this book in one way: their goal is set by a liability or a mandate from a government or a board, not by a profit to
owners. A pension fund's quants work to meet benefits; a sovereign fund's to meet a return target with a risk budget set
by its ministry; neither is paid a share of what it makes.

\begin{definition}[Internal management]\label{def:in:asset-managers-pensions-sovereign-funds-and-insurers:internal}
\emph{Internal management}\index{internal management} is the management of an \omterm{def:in:asset-managers-pensions-sovereign-funds-and-insurers:owner}{asset owner}'s assets by its own staff
rather than by external managers under mandates; its share is the fraction of assets managed that way.
\end{definition}

\section{Where quants sit: allocation, risk, liabilities, factors and indices}

In an \omterm{def:in:asset-managers-pensions-sovereign-funds-and-insurers:owner}{asset owner}, quantitative work falls into five kinds.
\begin{enumerate}
\item \textbf{Asset allocation}: the strategic mix and its risk, the owner's largest single decision.
\item \textbf{Risk}: measurement and limits across internal and external mandates (Book~6).
\item \textbf{Liability matching}: for pension funds and insurers, the liability-driven investment of Book~2 (chapter~7):
curves, duration, inflation, hedging programmes.
\item \textbf{Factor and quantitative equity}: internal systematic portfolios and the monitoring of external ones
(Book~8, chapter~29).
\item \textbf{Index and implementation}: low-cost replication of benchmarks, trading, transition management.
\end{enumerate}
The mix depends on the owner's choice between doing and buying. An owner that manages most assets internally employs
portfolio managers, traders and researchers; one that buys management employs manager selectors, risk analysts and
allocators.

\section{Two owners, two models}

\begin{dated}{2026-05}{Two asset owners in their own numbers}\label{dat:in:asset-managers-pensions-sovereign-funds-and-insurers:owners}
\textbf{Norges Bank Investment Management} (annual report 2025): fund value NOK 21\,268 billion; 678 employees in Oslo,
London, New York and Singapore; management fee NOK 7\,537 million, 3.8 basis points of assets under management, of which
personnel costs NOK 2\,476 million, base fees to external managers NOK 1\,756 million and performance fees NOK 1\,182
million; NOK 1\,062 billion (5.0\% of the fund) with 111 external mandates. \textbf{CPP Investments} (fiscal year to 31
March 2026): net assets C\$793.3 billion; operating expenses C\$1\,757 million (an operating expense ratio of 23.1 basis
points); C\$1\,976 million of investment management fees and C\$2\,758 million of performance fees paid to external
managers; C\$753 million of transaction costs; net investments per employee C\$364 million.
\end{dated}

The two numbers the owners publish, 3.8 and 23.1 basis points, are not comparable: one includes external managers' fees
and the other does not. Put on the same footing from the published totals, with year-end assets as the denominator, the
Norwegian fund's total management cost is 3.5 basis points and the Canadian fund's 81.8, of which 59.7 are external
managers' fees (\cref{fig:in:asset-managers-pensions-sovereign-funds-and-insurers:costs}). Splitting the Norwegian fund
by who manages the assets shows why: its internally managed 95\% cost 2.3 basis points, its external mandates 27.7
basis points, of which 16.5 base fees and the rest performance fees.

\begin{omfigure}
\begin{tikzpicture}
\begin{axis}[width=10.5cm,height=5.4cm,xbar,bar width=9pt,xmin=0,xmax=90,xlabel={\small management cost (basis points of assets)},
  ytick={0,1,2,3,4},yticklabels from table={figdata/industry/08-asset-managers-pensions-sovereign-funds-and-insurers/costs.csv}{label},
  y dir=reverse,tick label style={font=\footnotesize},grid=major,grid style={gray!20},table/col sep=comma,
  nodes near coords={\pgfmathprintnumber[fixed,precision=1]{\pgfplotspointmeta}},nodes near coords style={font=\scriptsize}]
\addplot[fill=omDef!55,draw=omDef] table[y=k,x=bp]{figdata/industry/08-asset-managers-pensions-sovereign-funds-and-insurers/costs.csv};
\end{axis}
\end{tikzpicture}

\omcaption{What two \omterm{def:in:asset-managers-pensions-sovereign-funds-and-insurers:owner}{asset owners} pay to manage their assets, in basis points of year-end assets, from their published
totals: the Norwegian fund's internally managed assets, its whole fund, the Canadian fund's own operating expenses, the
Norwegian fund's external mandates (base and performance fees), and the Canadian fund's total including external
managers' fees. Data: annual report 2025 and fiscal-2026 results release, through
\texttt{in\_owners}.}\label{fig:in:asset-managers-pensions-sovereign-funds-and-insurers:costs}
\end{omfigure}

\begin{remark}[What the two models buy]\label{rem:in:asset-managers-pensions-sovereign-funds-and-insurers:models}
A cost in basis points is half of the comparison: the other half is what the costs buy. An owner that pays external
managers large performance fees is paying for returns it expects them to earn above their fees, in assets (private
equity, credit, real assets, hedge funds) that an index team cannot run. An owner with 95\% of its assets in listed
securities managed close to an index needs few people per dollar. The comparison is between business models, not
between efficiencies.
\end{remark}

For an employee the models mean different jobs. The Norwegian fund manages \$3.0 billion of assets per employee, and
its personnel costs come to NOK 3.65 million per employee, about \$352\,000 at 2025 average rates: fewer people, each
responsible for more money, working mostly on listed markets. The Canadian fund's model employs more people per dollar
(C\$364 million of net investments per employee) and many of them select, oversee and negotiate with external managers
and partners, which is a different craft from running a portfolio.

\section{Build or hire: the internal team's break-even}

\begin{method}[Break-even mandate size]\label{met:in:asset-managers-pensions-sovereign-funds-and-insurers:be}
\begin{enumerate}
\item Cost the internal team: headcount times staff cost per head, plus systems and data.
\item Cost the external option: the fee rate in basis points of the mandate, with the expected performance fee.
\item The break-even mandate size is the internal cost divided by the fee rate.
\item Above it the internal team is cheaper, if it can deliver the same return; below it, the fee is.
\end{enumerate}
\end{method}

\begin{example}[A ten-person team]\label{ex:in:asset-managers-pensions-sovereign-funds-and-insurers:team}
A ten-person internal team at \$352\,000 per head, the Norwegian fund's 2025 personnel cost, plus \$1.5 million of systems
and data (an illustrative figure), costs \$5.02 million a year. Against an external fee of 20 basis points it breaks even at a
mandate of \$2.51 billion; against 10 basis points, at \$5.02 billion; against 40, at \$1.255 billion
(\cref{fig:in:asset-managers-pensions-sovereign-funds-and-insurers:be}). A large owner's mandates are far above these
sizes, which is why the largest owners internalise the strategies they can run.
\end{example}

\begin{omfigure}
\begin{tikzpicture}
\begin{axis}[width=10.5cm,height=5.8cm,xmin=5,xmax=60,ymin=0,ymax=20,xlabel={\small external fee (basis points of the mandate)},
  ylabel={\small break-even mandate (\$ billion)},tick label style={font=\footnotesize},grid=major,grid style={gray!20},
  table/col sep=comma,legend style={legend cell align=left,font=\scriptsize,at={(0.5,-0.3)},anchor=north,/tikz/every even column/.append style={column sep=8pt}},legend columns=3]
\addplot[omDef,thick] table[x=fee_bp,y=team5]{figdata/industry/08-asset-managers-pensions-sovereign-funds-and-insurers/breakeven.csv};
\addplot[omThm,thick,dashed] table[x=fee_bp,y=team10]{figdata/industry/08-asset-managers-pensions-sovereign-funds-and-insurers/breakeven.csv};
\addplot[omProp,thick,dotted] table[x=fee_bp,y=team20]{figdata/industry/08-asset-managers-pensions-sovereign-funds-and-insurers/breakeven.csv};
\legend{5 people,10 people,20 people}
\end{axis}
\end{tikzpicture}

\omcaption{The mandate size above which an internal team is cheaper than an external manager, against the external fee,
for teams of 5, 10 and 20 people at \$352\,000 per head plus \$1.5 million of systems and data (illustrative). Data:
\texttt{firm.ownercost.breakeven\_assets}, through \texttt{in\_owners.team\_breakeven}.}\label{fig:in:asset-managers-pensions-sovereign-funds-and-insurers:be}
\end{omfigure}

\section{Long-only quantitative managers}

The asset managers that serve owners are the largest employers on the buy side. One manager reports \$14.0 trillion of
assets under management and about 24\,900 employees, some \$560 million per employee. Quantitative work at a manager of
that size is index
engineering, portfolio construction and factor products (Book~8, chapter~29), risk systems and the trading that
implements portfolio changes (Book~10, chapter~20). The pace is set by rebalancing calendars and client mandates, not by
the market's minute-to-minute moves, and pay follows the manager's fee revenue rather than trading profit
(chapter~13).

\section{Insurers: asset-liability management and the actuarial quant}

An insurer invests premiums against liabilities it owes policyholders, often decades away; its investment function is
an \omterm{def:in:asset-managers-pensions-sovereign-funds-and-insurers:owner}{asset owner}'s, and its quants work on asset--liability management (Book~6, chapter~24): matching the duration and
cash flows of assets to liabilities, hedging interest-rate and inflation risk, and holding capital against the mismatch
under solvency rules. Beside them, actuaries price the liabilities themselves with mortality, lapse and
catastrophe models. The two professions meet on the balance sheet, and a quant moving between banking and insurance
brings interest-rate and derivatives models (Books~2 and~6) and learns liability models.

\section{Pace, pay and stability on the asset-owner side}

The asset-owner side offers a different bargain from the rest of this book. Horizons are long, feedback is slow and
measured against benchmarks, and a single decision --- the strategic allocation --- outweighs any trade. Pay is set by
boards and governments, published in aggregate (the personnel line of the Norwegian fund's accounts), and rarely tied
to a share of profit. Institutions last: a pension fund does not close after a drawdown. For a quant, the side trades
upside for stability and scale for pace.

\section{Tutorial: build or hire the team}\label{tut:in:asset-managers-pensions-sovereign-funds-and-insurers:tut}

\begin{tutorial}
\textbf{Goal.} Put two \omterm{def:in:asset-managers-pensions-sovereign-funds-and-insurers:owner}{asset owners}' published costs on one footing, split one owner's cost between internal and external
management, and compute the break-even mandate for an internal team. \textbf{End state:}
\cref{fig:in:asset-managers-pensions-sovereign-funds-and-insurers:costs,fig:in:asset-managers-pensions-sovereign-funds-and-insurers:be}.
\begin{tutsteps}
\item \textbf{The totals.} \texttt{data/industry/asset\_owners.csv} holds each owner's published totals in its own
currency: assets, employees, personnel and internal costs, external base and performance fees, externally managed
assets where published.
\item \textbf{One footing.} \texttt{firm.ownercost.cost\_bp(owner)} divides all management costs by year-end assets;
\texttt{split\_bp} divides internal costs by internally managed assets and external fees by externally managed assets
(\cref{lst:in:asset-managers-pensions-sovereign-funds-and-insurers:split}).
\omcode{code/firm/ownercost/firm_ownercost.py}{30}{38}{The total cost and the internal-external split, in basis points.}\label{lst:in:asset-managers-pensions-sovereign-funds-and-insurers:split}
\item \textbf{Per head.} \texttt{in\_owners.nbim()} converts personnel costs per employee to dollars at the ECB's average
rates.
\item \textbf{Break-even.} \texttt{in\_owners.team\_breakeven(heads, cost\_per\_head\_k, systems\_k, fees)} applies
\cref{met:in:asset-managers-pensions-sovereign-funds-and-insurers:be}.
\end{tutsteps}
\end{tutorial}

\textbf{What to change next.} Add a performance fee of 20\% of a 1\% excess return to the external option and recompute
the break-even (exercise~7); add a third owner from its annual report, with its own cost lines.

\section{Build: the owner's cost model}\label{bld:in:asset-managers-pensions-sovereign-funds-and-insurers:ownercost}

\begin{build}
\bfield{Purpose} Compare doing and buying investment management on one footing, for this chapter and for chapter~12's
per-head comparisons.
\bfield{Interface} \texttt{firm.ownercost}: \texttt{Owner(name, assets, internal\_costs, external\_base, external\_perf,
external\_assets)}; \texttt{cost\_bp}; \texttt{split\_bp}; \texttt{internal\_cost}; \texttt{external\_cost};
\texttt{breakeven\_assets}.
\bfield{Rules} One currency and one denominator per comparison; the split is refused when the owner does not publish its
externally managed assets; fees in basis points of the mandate.
\bfield{Acceptance tests} \texttt{code/firm/ownercost/tests/}: cost and split on constructed totals; the split refused
without external assets; internal and external costs equal at the break-even.
\bfield{Stretch} Transaction costs and the return difference between internal and external management; a team whose
cost per head rises with the mandate's complexity.
\end{build}

\begin{omsources}
\item Norges Bank Investment Management, Government Pension Fund Global, annual report 2025 (table 12.1).
\item CPP Investments, results release for fiscal 2026 (20 May 2026).
\item International Working Group of Sovereign Wealth Funds, Santiago Principles (2008); US Internal Revenue Service,
defined benefit plans.
\item BlackRock, Form 10-K for 2025; European Central Bank, annual average reference rates.
\end{omsources}

\section{Exercises}

\begin{exercise}[$\star$]\label{exo:in:asset-managers-pensions-sovereign-funds-and-insurers:1}
Convert the Norwegian fund's value to US dollars at the ECB's 2025 average rates, and compute its assets per employee.
\end{exercise}

\begin{exercise}[$\star$]\label{exo:in:asset-managers-pensions-sovereign-funds-and-insurers:2}
Compute the Norwegian fund's base fees to external managers in basis points of the externally managed assets.
\end{exercise}

\begin{exercise}[$\star$]\label{exo:in:asset-managers-pensions-sovereign-funds-and-insurers:3}
A pension plan promises a fixed benefit at retirement. Who bears the investment risk, and which quant work follows?
\end{exercise}

\begin{exercise}[$\star\star$]\label{exo:in:asset-managers-pensions-sovereign-funds-and-insurers:4}
Compute the Canadian fund's external managers' fees and its own operating expenses in basis points of year-end net
assets, and explain why its published operating expense ratio (23.1) differs from your operating figure.
\end{exercise}

\begin{exercise}[$\star\star$]\label{exo:in:asset-managers-pensions-sovereign-funds-and-insurers:5}
A 20-person internal team costs \$352\,000 per head plus \$1.5 million of systems. What mandate size breaks even against an
external fee of 25 basis points?
\end{exercise}

\begin{exercise}[$\star\star$]\label{exo:in:asset-managers-pensions-sovereign-funds-and-insurers:6}
Read \cref{fig:in:asset-managers-pensions-sovereign-funds-and-insurers:be}: for a ten-person team, what fee makes a
\$2 billion mandate break even?
\end{exercise}

\begin{exercise}[$\star\star\star$]\label{exo:in:asset-managers-pensions-sovereign-funds-and-insurers:7}
\emph{Coding.} An external manager charges 20 basis points plus 20\% of an expected excess return of 1\% a year. Recompute
the ten-person team's break-even mandate. What must the internal team deliver for the comparison to hold?
\end{exercise}

\begin{exercise}[$\star\star\star$]\label{exo:in:asset-managers-pensions-sovereign-funds-and-insurers:8}
\emph{Find the flaw.} ``The Norwegian fund costs 3.5 basis points and the Canadian fund 81.8: the Canadian fund is 23 times
less efficient.''
\end{exercise}

\section{Problem: Build or Hire the Team}

\begin{problem}[{Weekend problem --- build or hire the team}]\label{pb:in:asset-managers-pensions-sovereign-funds-and-insurers:1}
The chief investment officer of a pension fund with \$200 billion of assets asks whether to bring a \$15 billion
equity mandate in-house, and a quant is asked to answer from public numbers.

\medskip\noindent\textbf{Part I --- The owners.}
\begin{enumerate}
\item Define an \omterm{def:in:asset-managers-pensions-sovereign-funds-and-insurers:owner}{asset owner}, a \omterm{def:in:asset-managers-pensions-sovereign-funds-and-insurers:swf}{sovereign wealth fund} and a \omterm{def:in:asset-managers-pensions-sovereign-funds-and-insurers:swf}{defined-benefit pension plan}.
\item Define \omterm{def:in:asset-managers-pensions-sovereign-funds-and-insurers:internal}{internal management} and give the Norwegian fund's internal share.
\item Name the five kinds of quantitative work in an \omterm{def:in:asset-managers-pensions-sovereign-funds-and-insurers:owner}{asset owner}.
\item Give the two owners' published cost ratios, and say why they cannot be compared.
\item What does the Norwegian fund's management fee consist of?
\end{enumerate}

\medskip\noindent\textbf{Part II --- One footing.}
\begin{enumerate}[resume]
\item Compute both owners' total management costs in basis points of year-end assets.
\item Split the Norwegian fund's cost between internal and external management.
\item What share of the Canadian fund's total cost is external managers' fees?
\item Give the Norwegian fund's personnel cost per employee in kroner and in dollars.
\item Give its assets per employee in dollars.
\end{enumerate}

\medskip\noindent\textbf{Part III --- The team.}
\begin{enumerate}[resume]
\item State the break-even method.
\item Cost a ten-person team at the Norwegian fund's personnel cost per head and \$1.5 million of systems.
\item Give its break-even mandate at external fees of 10, 20 and 40 basis points.
\item At 20 basis points, how many times the break-even is the \$15 billion mandate?
\item What would the external manager cost on \$15 billion at 20 basis points?
\end{enumerate}

\medskip\noindent\textbf{Part IV --- The verdict.}
\begin{enumerate}[resume]
\item State the \emph{named result}: the break-even mandate sizes for a ten-person team at 10, 20 and 40 basis points.
\item What must the internal team deliver for the answer to hold?
\item What does the comparison of the two owners' costs leave out?
\item Which jobs does bringing the mandate in-house create, and which does it remove?
\item In two sentences, what should the quant tell the chief investment officer?
\end{enumerate}
\end{problem}

\section{Interview questions}

\begin{interviewq}[$\star$ \iqroles{researcher}]\label{iq:in:asset-managers-pensions-sovereign-funds-and-insurers:1}
What is the difference between an \omterm{def:in:asset-managers-pensions-sovereign-funds-and-insurers:owner}{asset owner} and an asset manager, and how does it change a quant's job?
\end{interviewq}

\begin{interviewq}[$\star$ \iqroles{risk}]\label{iq:in:asset-managers-pensions-sovereign-funds-and-insurers:2}
A pension plan's liabilities have a duration of 18 years and its bonds a duration of 6. What happens to its funding
ratio when long rates fall, and what would you do about it?
\end{interviewq}

\begin{interviewq}[$\star\star$ \iqroles{researcher}]\label{iq:in:asset-managers-pensions-sovereign-funds-and-insurers:3}
An owner pays an external manager 50 basis points plus 20\% of outperformance. What excess return must the manager
expect to earn to be worth hiring over a 5-basis-point internal index portfolio?
\end{interviewq}

\begin{interviewq}[$\star\star$ \iqroles{researcher, risk}]\label{iq:in:asset-managers-pensions-sovereign-funds-and-insurers:4}
Why does the strategic asset allocation dominate a large fund's return, and what does that imply for where its best
quants should work?
\end{interviewq}

\begin{interviewq}[$\star\star$ \iqroles{bank, risk}]\label{iq:in:asset-managers-pensions-sovereign-funds-and-insurers:5}
Which of your skills as a bank rates quant transfer to an insurer's asset--liability team, and what would you have to
learn?
\end{interviewq}

\begin{interviewq}[$\star\star\star$ \iqroles{researcher}]\label{iq:in:asset-managers-pensions-sovereign-funds-and-insurers:6}
Design a test of whether an owner's internal equity team adds value net of its costs, given ten years of monthly returns
of the team and its benchmark.
\end{interviewq}
